%
% connection_lease_reaping_cascade_incomplete_buggy.tex
%
% FIDELITY CONTROL for round 3's I7 (cascade-incomplete), not the design.
% Changes only TimedReapEnd and Claim: both omit the three new (+) terms
% this round adds, leaving hasSubs/hasWriter/inboxNonEmpty exactly as they
% were instead of clearing them to fclear -- the literal shape of the gap
% docs/architecture/lease-reap-cascade.md Section 1 reports: Depart, once
% fixed, already reaches all four cascade legs via lifecycle.py, but
% _depart_lapsed (TimedReap's realization) and apply's embedded sweep
% (Claim) never called Hub.on_disconnect or the transport sink at all.
% See docs/connection_lease_reaping.tex, Section "Cascade-incomplete", for
% the full trace and the goals that must return FOUND against this
% variant and NOT found against the fixed spec.
%
% Type-check:  fuzz docs/connection_lease_reaping_cascade_incomplete_buggy.tex
% Model-check:
%   probcli docs/connection_lease_reaping_cascade_incomplete_buggy.tex \
%       -model_check -nodead -goal "..." -p DEFAULT_SETSIZE 3
%

\documentclass[a4paper,10pt,fleqn]{article}
\usepackage[margin=1in]{geometry}
\usepackage{fuzz}
\usepackage[colorlinks=true,linkcolor=blue,citecolor=blue,urlcolor=blue]{hyperref}

\begin{document}

\title{lux Connection-Lease Reaping --- Fidelity Control (Cascade Incomplete)}
\author{Buggy variant: TimedReapEnd and Claim release the registry and
  ownership but never clear subscriptions, the writer binding, or the
  inbox}
\date{September 2026}
\maketitle

% ============================================================
\section{Introduction}
% ============================================================

This is the fidelity control for round 3's I7 in
\texttt{docs/connection\_lease\_reaping.tex}, not the design. It removes
only the three new cascade-clearing terms \texttt{TimedReapEnd} and
\texttt{Claim} otherwise carry: the registry-and-ownership legs are
released exactly as the fixed spec releases them, but \texttt{hasSubs},
\texttt{hasWriter}, and \texttt{inboxNonEmpty} are left untouched.
\texttt{\_depart\_lapsed} and \texttt{apply}'s embedded sweep never call
\texttt{Hub.on\_disconnect} or the transport-owned sink today -- a
timer-reaped or swept-aside connection keeps its subscriptions, writer
binding, and inbox queue forever, even though the identical connection is
already gone from the registry and owns nothing. Read the main
specification first; this document states only what differs, in
Section~\ref{sec:timedreapend} and Section~\ref{sec:claim}.

% ============================================================
\section{Basic Types}
% ============================================================

\begin{zed}
  [CONN, IDENT, SCENE]
\end{zed}

\begin{zed}
  TSTATE ::= tup | tdown
\end{zed}

\begin{zed}
  LSTATE ::= llive | llapsed
\end{zed}

\begin{zed}
  OWNER ::= unowned | conn \ldata CONN \rdata
\end{zed}

\begin{zed}
  FSTATE ::= fclear | fset
\end{zed}

Unchanged from the main specification.

% ============================================================
\section{State}
% ============================================================

\begin{schema}{ConnReg}
  registered : \power CONN \\
  identity : CONN \pfun IDENT \\
  transport : CONN \pfun TSTATE \\
  lease : CONN \pfun LSTATE \\
  owner : SCENE \fun OWNER \\
  hasSubs : CONN \fun FSTATE \\
  hasWriter : CONN \fun FSTATE \\
  inboxNonEmpty : CONN \fun FSTATE \\
  cascadeFor : \power CONN
\where
  registered \subseteq \dom identity \\
  \dom identity = \dom transport \\
  \dom transport = \dom lease \\
  \# cascadeFor \leq 1
\end{schema}

Unchanged from the main specification.

% ============================================================
\section{Initialization}
% ============================================================

\begin{schema}{Init}
  ConnReg'
\where
  registered' = \emptyset \\
  identity' = \emptyset \\
  transport' = \emptyset \\
  lease' = \emptyset \\
  owner' = (\lambda s : SCENE @ unowned) \\
  hasSubs' = (\lambda c : CONN @ fclear) \\
  hasWriter' = (\lambda c : CONN @ fclear) \\
  inboxNonEmpty' = (\lambda c : CONN @ fclear) \\
  cascadeFor' = \emptyset
\end{schema}

Unchanged from the main specification.

% ============================================================
\section{Operations}
% ============================================================

\begin{schema}{Connect}
  \Delta ConnReg \\
  c? : CONN \\
  i? : IDENT
\where
  cascadeFor = \emptyset \\
  registered' = registered \cup \{ c? \} \\
  identity' = identity \oplus \{ c? \mapsto i? \} \\
  transport' = transport \oplus \{ c? \mapsto tup \} \\
  lease' = lease \oplus \{ c? \mapsto llive \} \\
  hasWriter' = hasWriter \oplus \{ c? \mapsto fset \} \\
  hasSubs' = hasSubs \\
  inboxNonEmpty' = inboxNonEmpty \\
  cascadeFor' = cascadeFor \\
  owner' = owner
\end{schema}

\begin{schema}{Renew}
  \Delta ConnReg \\
  c? : CONN
\where
  c? \in registered \\
  transport~c? = tup \\
  cascadeFor = \emptyset \\
  lease' = lease \oplus \{ c? \mapsto llive \} \\
  registered' = registered \\
  identity' = identity \\
  transport' = transport \\
  hasSubs' = hasSubs \\
  hasWriter' = hasWriter \\
  inboxNonEmpty' = inboxNonEmpty \\
  cascadeFor' = cascadeFor \\
  owner' = owner
\end{schema}

\begin{schema}{Activate}
  \Delta ConnReg \\
  c? : CONN
\where
  c? \in registered \\
  transport~c? = tup \\
  hasSubs' = hasSubs \oplus \{ c? \mapsto fset \} \\
  hasWriter' = hasWriter \oplus \{ c? \mapsto fset \} \\
  inboxNonEmpty' = inboxNonEmpty \oplus \{ c? \mapsto fset \} \\
  registered' = registered \\
  identity' = identity \\
  transport' = transport \\
  lease' = lease \\
  owner' = owner \\
  cascadeFor' = cascadeFor
\end{schema}

\begin{schema}{TransportDies}
  \Delta ConnReg \\
  c? : CONN
\where
  c? \in registered \\
  transport~c? = tup \\
  transport' = transport \oplus \{ c? \mapsto tdown \} \\
  registered' = registered \\
  identity' = identity \\
  lease' = lease \\
  owner' = owner \\
  hasSubs' = hasSubs \\
  hasWriter' = hasWriter \\
  inboxNonEmpty' = inboxNonEmpty \\
  cascadeFor' = cascadeFor
\end{schema}

\begin{schema}{LeaseLapse}
  \Delta ConnReg \\
  c? : CONN
\where
  c? \in registered \\
  lease~c? = llive \\
  lease' = lease \oplus \{ c? \mapsto llapsed \} \\
  registered' = registered \\
  identity' = identity \\
  transport' = transport \\
  owner' = owner \\
  hasSubs' = hasSubs \\
  hasWriter' = hasWriter \\
  inboxNonEmpty' = inboxNonEmpty \\
  cascadeFor' = cascadeFor
\end{schema}

\begin{schema}{Read}
  \Xi ConnReg
\end{schema}

\begin{schema}{Depart}
  \Delta ConnReg \\
  c? : CONN
\where
  c? \in registered \\
  registered' = registered \setminus \{ c? \} \\
  owner' = owner \oplus \\
  \quad~ (\{ s : SCENE | owner~s = conn(c?) \} \cross \{ unowned \}) \\
  hasSubs' = hasSubs \oplus \{ c? \mapsto fclear \} \\
  hasWriter' = hasWriter \oplus \{ c? \mapsto fclear \} \\
  inboxNonEmpty' = inboxNonEmpty \oplus \{ c? \mapsto fclear \} \\
  identity' = identity \\
  transport' = transport \\
  lease' = lease \\
  cascadeFor' = cascadeFor
\end{schema}

\texttt{Depart} is unchanged from the fixed spec, and correct: this
control isolates the gap to \texttt{TimedReapEnd} and \texttt{Claim}
alone, matching the design document's own report that \texttt{Depart}'s
real-code counterpart already reaches all four cascade legs once fixed.

\begin{schema}{TimedReapBegin}
  \Delta ConnReg \\
  c? : CONN
\where
  c? \in registered \\
  lease~c? = llapsed \\
  cascadeFor = \emptyset \\
  registered' = registered \setminus \{ c? \} \\
  owner' = owner \oplus \\
  \quad~ (\{ s : SCENE | owner~s = conn(c?) \} \cross \{ unowned \}) \\
  cascadeFor' = \{ c? \} \\
  identity' = identity \\
  transport' = transport \\
  lease' = lease \\
  hasSubs' = hasSubs \\
  hasWriter' = hasWriter \\
  inboxNonEmpty' = inboxNonEmpty
\end{schema}

Unchanged from the fixed spec.

\subsection{TimedReapEnd --- BUGGY: releases the lock, never the cascade}
\label{sec:timedreapend}

\begin{schema}{TimedReapEnd}
  \Delta ConnReg \\
  c? : CONN
\where
  cascadeFor = \{ c? \} \\
  cascadeFor' = \emptyset \\
  hasSubs' = hasSubs \\
  hasWriter' = hasWriter \\
  inboxNonEmpty' = inboxNonEmpty \\
  registered' = registered \\
  identity' = identity \\
  transport' = transport \\
  lease' = lease \\
  owner' = owner
\end{schema}

The three terms that clear \texttt{hasSubs}/\texttt{hasWriter}/
\texttt{inboxNonEmpty} in the fixed spec's \texttt{TimedReapEnd} are
replaced here with the identity -- the cascade window closes
(\texttt{cascadeFor} still clears), but nothing the window was meant to
release actually is.

\subsection{Claim --- BUGGY: sweeps the registry and ownership, never the
cascade}
\label{sec:claim}

\begin{schema}{Claim}
  \Delta ConnReg \\
  c? : CONN \\
  s? : SCENE
\where
  c? \in registered \\
  transport~c? = tup \\
  lease' = lease \oplus \{ c? \mapsto llive \} \\
  registered' = registered \setminus \\
  \quad~ \{ c2 : CONN | c2 \in registered \land c2 \neq c? \land
    lease~c2 = llapsed \} \\
  (owner \oplus (\{ s : SCENE | \exists c2 : CONN @ \\
    \qquad c2 \in registered \land c2 \neq c? \land lease~c2 = llapsed
    \land owner~s = conn(c2) \} \\
    \qquad \cross \{ unowned \}))~s? = unowned \\
  \lor \\
  (owner \oplus (\{ s : SCENE | \exists c2 : CONN @ \\
    \qquad c2 \in registered \land c2 \neq c? \land lease~c2 = llapsed
    \land owner~s = conn(c2) \} \\
    \qquad \cross \{ unowned \}))~s? = conn(c?) \\
  owner' = (owner \oplus (\{ s : SCENE | \exists c2 : CONN @ \\
    \qquad c2 \in registered \land c2 \neq c? \land lease~c2 = llapsed
    \land owner~s = conn(c2) \} \\
    \qquad \cross \{ unowned \})) \oplus \{ s? \mapsto conn(c?) \} \\
  hasSubs' = hasSubs \\
  hasWriter' = hasWriter \\
  inboxNonEmpty' = inboxNonEmpty \\
  identity' = identity \\
  transport' = transport \\
  cascadeFor' = cascadeFor
\end{schema}

The registry-and-ownership sweep is exactly the fixed spec's: same set,
same \(c2 \neq c?\) exclusion. Only the three cascade-clearing terms are
replaced with the identity.

% ============================================================
\section{Verification: the goal that must flip}
% ============================================================

Against this variant, I7's negation must be \emph{found} -- the trace in
\texttt{connection\_lease\_reaping.tex}, Section ``Cascade-incomplete,''
reaches it in six steps:
\texttt{Connect}(c1, i1); \texttt{Activate}(c1); \texttt{TransportDies}(c1);
\texttt{LeaseLapse}(c1); \texttt{TimedReapBegin}(c1); \texttt{TimedReapEnd}(c1).

\begin{verbatim}
  # must be FOUND against this variant (NOT found against the fixed spec)
  probcli docs/connection_lease_reaping_cascade_incomplete_buggy.tex \
    -model_check -nodead \
    -goal "#c.(c:CONN & not(c : registered) & not(c : cascadeFor) &
               (hasSubs(c) = fset or hasWriter(c) = fset or
                inboxNonEmpty(c) = fset))" \
    -p DEFAULT_SETSIZE 3
\end{verbatim}

The identical goal is also reachable through \texttt{Claim}'s half of the
same gap, with no \texttt{TimedReapBegin}/\texttt{TimedReapEnd} in the
trace at all:
\texttt{Connect}(c1, i1); \texttt{Activate}(c1); \texttt{LeaseLapse}(c1);
\texttt{Connect}(c2, i2); \texttt{Claim}(c2, s2) -- the last step's sweep
removes \texttt{c1} and releases its ownership, but leaves its three
cascade flags untouched.

I2/I3/I5/I6 must stay unaffected -- this variant does not touch
\texttt{registered} or \texttt{owner} differently from the fixed spec
anywhere:

\begin{verbatim}
  # must remain NOT found
  probcli docs/connection_lease_reaping_cascade_incomplete_buggy.tex \
    -model_check -nodead \
    -goal "#s.(#c.(s:SCENE & c:CONN & owner(s) = conn(c) &
                   not(c : registered)))" \
    -p DEFAULT_SETSIZE 3
\end{verbatim}

I1 must stay unaffected -- this variant does not touch
\texttt{TransportDies}, \texttt{LeaseLapse}, or \texttt{TimedReapBegin},
so the chain that establishes I1's positive outcome still completes:

\begin{verbatim}
  # must remain FOUND
  probcli docs/connection_lease_reaping_cascade_incomplete_buggy.tex \
    -model_check -nodead \
    -goal "#c.(c:CONN & c : dom(identity) & not(c : registered) &
               transport(c) = tdown)" \
    -p DEFAULT_SETSIZE 3
\end{verbatim}

\end{document}
